% Part: model-theory
% Chapter: interpolation
% Section: definability

\documentclass[../../../include/open-logic-section]{subfiles}

\begin{document}

\olfileid{mod}{int}{def}

\olsection{قضیهٔ تعریف‌پذیری}

یکی از کاربردهای مهم قضیهٔ درون‌یابی، قضیهٔ تعریف‌پذیری بث است. برای
تعریف یک رابطهٔ $n$موضعی~$R$ می‌توانیم !!a{formula}~$!C$ با $n$
!!{variable}s آزاد بدهیم که~$R$ در آن دخیل نیست. این تعریفی
\emph{صریح} از~$R$ بر حسب~$!C$ خواهد بود. سپس می‌توانیم بگوییم
نظریهٔ~$\Sigma(P)$ در !!{language}ی شامل !!{predicate}~$P$
$n$موضعی، $P$ را صریحاً تعریف می‌کند اگر تعریفی صریح و صوری‌شده را در
بر داشته باشد (یا دست‌کم نتیجه دهد)؛ یعنی
\[
\Sigma(P) \Entails \lforall[x_1][\dots
  \lforall[x_n][(\Atom{P}{x_1,\dots, x_n} \liff !C(x_1, \dots,
    x_n))]].
\]
اما تعریف صریح تنها یک راه برای تعریف‌کردن---به معنای تعیین کامل---یک
رابطه است. همچنین ممکن است نظریه‌ای چنان باشد که در هر مدل، تعبیر~$P$
با تعبیر باقی !!{language} تثبیت شود. قضیهٔ تعریف‌پذیری می‌گوید هرگاه
نظریه‌ای تعبیر~$P$ را به این شیوه تثبیت کند---هرگاه~$P$ را
\emph{به‌طور ضمنی تعریف کند}---آن را به‌طور صریح نیز تعریف می‌کند.

\begin{defn}
فرض کنید $\Lang{L}$ !!a{language} باشد که !!{predicate}~$P$ را در بر
ندارد. مجموعهٔ~$\Sigma(P)$ از !!{sentence}s~$\Lang{L} \cup \{P\}$،
$P$ را \emph{صریحاً تعریف می‌کند} هرگاه و تنها هرگاه
!!a{formula}~$!C(x_1, \dots, x_n)$ از~$\Lang{L}$ وجود داشته باشد
چنان‌که
\[
\Sigma(P) \Entails \lforall[x_1][\dots
  \lforall[x_n][(\Atom{P}{x_1,\dots, x_n} \liff !C(x_1, \dots,
    x_n))]].
\]
\end{defn}

\begin{defn}
فرض کنید $\Lang{L}$ !!a{language} باشد که !!{predicate}s~$P$ و~$P'$
را در بر ندارد. مجموعهٔ~$\Sigma(P)$ از !!{sentence}s
$\Lang{L} \cup \{P\}$، $P$ را \emph{به‌طور ضمنی تعریف می‌کند}
هرگاه و تنها هرگاه
\[
\Sigma(P) \cup \Sigma(P') \Entails \lforall[x_1][\dots
  \lforall[x_n][(\Atom{P}{x_1,\dots, x_n} \liff \Atom{P'}{x_1,\dots,
      x_n})]],
\]
که در آن $\Sigma(P')$ حاصل جایگزینی یکنواخت~$P$ با~$P'$ در
$\Sigma(P)$ است.
\end{defn}

به بیان دیگر، برای هر مدل~$\Struct{M}$ و
$R, R' \subseteq \Domain{M}^n$، اگر هم
$\Sat{\Expan{M}{R}}{\Sigma(P)}$ و هم
$\Sat{\Expan{M}{R'}}{\Sigma(P')}$، آنگاه $R=R'$؛ که در آن
$\Expan{M}{R}$ همان !!{structure}~$\Struct{M'}$ برای بسط~$\Lang{L}$
به~$\Lang{L} \cup \{P\}$ است چنان‌که $\Assign{P}{M'} = R$، و برای
$\Expan{M}{R'}$ نیز به همین ترتیب.

\begin{thm}[قضیهٔ تعریف‌پذیری بث] مجموعهٔ~$\Sigma(P)$ از
  !!{sentence}s-$\Lang{L} \cup\{P\}$، $P$ را به‌طور ضمنی تعریف می‌کند
  اگر و تنها اگر $\Sigma(P)$، $P$ را به‌طور صریح تعریف کند.
\end{thm}

\begin{proof}
اگر $\Sigma(P)$، $P$ را به‌طور صریح تعریف کند، آنگاه هر دو گزارهٔ
\begin{align*}
  \Sigma(P) & \Entails & \lforall[x_1][\dots \lforall[x_n]
    [(\Atom{P}{x_1,\dots, x_n} \liff !C(x_1,\dots,x_n))]]\\
  \Sigma(P') & \Entails & \lforall[x_1][\dots \lforall[x_n]
    [(\Atom{P'}{x_1,\dots, x_n} \liff !C(x_1,\dots,x_n))]]
\end{align*}
برقرارند و نتیجه به دست می‌آید. برای عکس، فرض کنید $\Sigma(P)$، $P$
را به‌طور ضمنی تعریف می‌کند. نخست !!{constant}s~$c_1$، \dots،~$c_n$
را به~$\Lang{L}$ می‌افزاییم. آنگاه
\[
\Sigma(P) \cup \Sigma(P') \Entails
\Atom{P}{c_1, \dots, c_n} \to  \Atom{P'}{c_1, \dots, c_n}.
\]
بنا بر فشردگی، مجموعه‌های متناهی
$\Delta_0 \subseteq \Sigma(P)$ و $\Delta_1 \subseteq \Sigma(P')$
وجود دارند چنان‌که
\[
\Delta_0 \cup \Delta_1 \Entails
\Atom{P}{c_1, \dots, c_n} \to \Atom{P'}{c_1, \dots, c_n}.
\]
بگذارید $!D(P)$ عطف همهٔ !!{sentence}s~$!A(P)$ باشد چنان‌که یا
$!A(P) \in \Delta_0$ یا $!A(P') \in \Delta_1$؛ و بگذارید $!D(P')$
عطف همهٔ !!{sentence}s~$!A(P')$ باشد چنان‌که یا
$!A(P) \in \Delta_0$ یا $!A(P') \in \Delta_1$. آنگاه
$!D(P) \land !D(P') \Entails \Atom{P}{c_1, \dots, c_n} \to
\Atom{P'}{c_1,\dots,c_n}$. می‌توانیم این عبارت را چنان بازآرایی کنیم
که هر !!{predicate} در یک سوی~$\Entails$ رخ دهد:
\[
!D(P) \land \Atom{P}{c_1, \dots, c_n} \Entails
!D(P') \to \Atom{P'}{c_1, \dots, c_n}.
\]
بنا بر قضیهٔ درون‌یابی کریگ، !!a{sentence}~$!C(c_1,\dots, c_n)$ که
شامل~$P$ یا~$P'$ نیست وجود دارد چنان‌که:
\begin{align*}
  !D(P) \land \Atom{P}{c_1, \dots, c_n} & \Entails !C(c_1,\dots, c_n); \\
  !C(c_1,\dots, c_n) & \Entails !D(P') \to \Atom{P'}{c_1, \dots, c_n}.
\end{align*}
از استلزام نخست داریم
$!D(P) \Entails \Atom{P}{c_1,\dots, c_n} \lif !C(c_1,\dots, c_n)$.
و از استلزام دوم، چون یک مدل~$\Lang{L} \cup \{P\}$ چون
$\Expan{M}{R}$ فرمول~$!A(P)$ را هرگاه و تنها هرگاه ارضا می‌کند که مدل
متناظر~$\Lang{L} \cup \{P'\}$ چون~$\Expan{M}{R}$، $!A(P')$ را ارضا
کند، داریم
$!C(c_1,\dots, c_n) \Entails !D(P) \lif \Atom{P}{c_1,\dots, c_n}$؛
و از این:
\[
!D(P) \Entails !C(c_1,\dots,c_n) \to \Atom{P}{c_1,\dots, c_n}.
\]
با ترکیب این دو،
$!D(P) \Entails \Atom{P}{c_1,\dots, c_n}
\liff !C(c_1, \dots, c_n)$؛ و بنا بر یکنوایی و تعمیم نیز
\[
\Sigma(P) \Entails
\lforall[x_1][\dots\lforall[x_n][(\Atom{P}{x_1,\dots, x_n} \liff
    !C(x_1,\dots, x_n))]].
\]
\end{proof}

\end{document}
